% ============================================================================
% CAPÍTULO 10: MODOS PARES E IMPARES
% ============================================================================
\chapter{Modos Pares e Impares: Lo que Klein Permite}
\label{cap:modos}

\begin{center}
\textit{``Las reglas de selección del universo''}
\end{center}

\vspace{0.5cm}

\begin{tcolorbox}[colback=white, colframe=kleinblue, boxrule=2pt]
\centering
Este capítulo explica uno de los aspectos más profundos de la Teoría Klein:
la selección de modos. La topología de la botella de Klein impone reglas
estrictas sobre qué configuraciones físicas son permitidas y cuáles no.
\end{tcolorbox}


% ============================================================================
\section{Nivel Narrativo: La Cinta de Möbius y el Espejo}
% ============================================================================

\begin{analogiabox}
\textbf{El Cañón Mágico}

Imagina un cañón mágico donde el eco vuelve invertido (como un espejo).
Si gritas ``¡HOLA!'', el eco devuelve ``¡ALOH!'' (invertido).

\textbf{MODOS PARES (permitidos):} Sonidos que suenan igual invertidos.

\begin{center}
``ANA'' $\to$ eco ``ANA'' $\checkmark$ (palíndromo, se refuerza)
\end{center}

\textbf{MODOS IMPARES (prohibidos):} Sonidos que se cancelan al invertirse.

\begin{center}
``HOLA'' $\to$ eco ``ALOH'' $\times$ (se interfieren destructivamente)
\end{center}

La botella de Klein es como este cañón: solo permite ``palabras palíndromo''
en el universo físico.
\end{analogiabox}

\textbf{¿POR QUÉ 7 CAPAS?}

La botella de Klein en física tiene 7 ``niveles de inversión'':
\begin{enumerate}
    \item Capa 1: Inversión espacial ($x \to -x$)
    \item Capa 2: Inversión temporal ($t \to -t$)
    \item Capa 3: Conjugación de carga ($q \to -q$)
    \item Capa 4: Inversión de helicidad
    \item Capa 5: Inversión de sabor
    \item Capa 6: Inversión de color (QCD)
    \item Capa 7: Inversión de generación
\end{enumerate}

Cada capa contribuye un factor $\pi$ al supresor total: $7\pi$.


% ============================================================================
\section{Nivel Conceptual: Funciones en Espacios No Orientables}
% ============================================================================

\textbf{CONDICIONES DE FRONTERA EN LA BOTELLA DE KLEIN:}

Una función $f$ definida sobre la botella de Klein debe satisfacer:

\begin{equation}
f(x, y) = f(x + L, -y) \quad \text{[Condición de Klein]}
\end{equation}

Esto significa que al dar una vuelta en $x$, la coordenada $y$ se invierte.

\textbf{DESCOMPOSICIÓN EN MODOS:}

Cualquier función puede escribirse como suma de modos:

\begin{equation}
f(x, y) = \sum_n \left[a_n \cos\left(\frac{n\pi y}{L}\right) + b_n \sin\left(\frac{n\pi y}{L}\right)\right] \times g(x)
\end{equation}

La condición de Klein impone:

\begin{itemize}
    \item Para modos PARES (coseno): $f(x, -y) = f(x, y)$ $\checkmark$ \textbf{PERMITIDO}
    \item Para modos IMPARES (seno): $f(x, -y) = -f(x, y)$ $\times$ \textbf{SUPRIMIDO}
\end{itemize}

\begin{tcolorbox}[colback=kleingray, colframe=kleinblue, title=\textbf{Regla de Selección de Klein}]
\begin{center}
\textbf{Amplitud física} $= A_0 \times (7\pi)^{-n}$

donde $n$ = número de inversiones requeridas por el modo
\end{center}
\end{tcolorbox}


% ============================================================================
\section{Tabla de Modos y Supresión}
% ============================================================================

\begin{table}[H]
\centering
\small
\caption{Cantidad física, paridad y supresión Klein}
\begin{tabular}{lccc}
\toprule
\textbf{Cantidad Física} & \textbf{Paridad} & \textbf{Supresión} & \textbf{Ejemplo} \\
\midrule
Velocidad de la luz & Par-Par & $(7\pi)^{-2}$ & $c = 3 \times 10^8 - \delta$ \\
Masa del electrón & Par & $(7\pi)^0$ & Escala fundamental \\
Masa del protón & Par$\times$Par & $6\pi^5$ & $m_p/m_e = 6\pi^5$ \\
Masa del Higgs & Par+Impar & $(7\pi)^1$ & Corrección al vacío \\
Const. estructura fina & Impar & $(7\pi)^1$ & $1/\alpha \approx 7^2\pi$ \\
Asimetría bariónica & Impar$\times$7 & $(7\pi)^{-7}$ & $\eta_B$ muy pequeño \\
Masa neutrino & Impar$\times$5 & $(7\pi)^{-5}$ & $m_\nu$ muy pequeña \\
Const. cosmológica & Impar$\times$92 & $(7\pi)^{-92}$ & $\Lambda$ extremadamente pequeña \\
\bottomrule
\end{tabular}
\end{table}


% ============================================================================
\section{Nivel Matemático: Teoría de Representaciones}
% ============================================================================

\begin{nivelmatematico}
\textbf{Estructura Matemática}

La botella de Klein $K^2$ tiene grupo fundamental:

\begin{equation}
\pi_1(K^2) = \langle a, b \,|\, aba^{-1}b = 1 \rangle
\end{equation}

Las representaciones de este grupo determinan qué campos son permitidos.

\textbf{Campos Escalares:}

Un campo escalar $\phi$ en $K^2$ debe satisfacer:

\begin{align}
\phi(T_a \cdot p) &= \chi_a \cdot \phi(p)\\
\phi(T_b \cdot p) &= \chi_b \cdot \phi(p)
\end{align}

donde $T_a$, $T_b$ son las traslaciones generadoras y $\chi$ son caracteres.

La relación del grupo fundamental impone:

\begin{equation}
\chi_a \cdot \chi_b \cdot \chi_a^{-1} \cdot \chi_b = 1
\end{equation}

Para $\chi \in U(1)$: $\chi_a^2 = 1$, por lo tanto $\chi_a = \pm 1$
\end{nivelmatematico}

\textbf{CLASIFICACIÓN DE MODOS:}

\begin{center}
\begin{tabular}{ccc}
\toprule
\textbf{Modo} $(\chi_a, \chi_b)$ & \textbf{Paridad} & \textbf{Permitido en Klein} \\
\midrule
$(+1, +1)$ & Par-Par & $\checkmark$ Completamente \\
$(+1, -1)$ & Par-Impar & $\checkmark$ Con supresión $(7\pi)^{-1}$ \\
$(-1, +1)$ & Impar-Par & $\checkmark$ Con supresión $(7\pi)^{-1}$ \\
$(-1, -1)$ & Impar-Impar & $\checkmark$ Con supresión $(7\pi)^{-2}$ \\
\bottomrule
\end{tabular}
\end{center}


% ============================================================================
\section{Aplicaciones Físicas Específicas}
% ============================================================================

\subsection{Velocidad de la Luz --- Modo Casi Puro}

El fotón es un modo casi completamente par. La pequeña impureza impar
produce la corrección:

\begin{equation}
c = c_0 \times \left(1 - \frac{1}{(7\pi)^2}\right)
\end{equation}

donde $c_0 = 3 \times 10^8$ m/s es el valor ``clásico'' par.

La corrección $1/(7\pi)^2 \approx 0.00207$ representa la contaminación del modo
impar al modo fundamental del fotón.

\subsection{Constante de Estructura Fina --- Modo Mixto}

La interacción electromagnética mezcla modos pares e impares:

\begin{align}
\frac{1}{\alpha} &= 7^2\pi - 7 - \pi^2\\
&= \underbrace{7^2\pi}_{\text{modo par principal}} - \underbrace{7}_{\text{corrección impar orden 1}} - \underbrace{\pi^2}_{\text{corrección orden 2}}
\end{align}

\subsection{Asimetría Bariónica --- 7 Inversiones}

La asimetría materia-antimateria requiere violar:
\begin{itemize}
    \item C (conjugación de carga)
    \item P (paridad)
    \item CP (combinación)
    \item T (reversión temporal, por CPT)
    \item B (número bariónico)
    \item Equilibrio térmico
    \item Simetría de sabor
\end{itemize}

Son exactamente 7 violaciones, dando:

\begin{equation}
\eta_B = \frac{3}{2} \times (7\pi)^{-7}
\end{equation}

\subsection{Constante Cosmológica --- Supresión Extrema}

La constante cosmológica es el modo MÁS impar posible. Requiere
cancelación en TODAS las escalas:

\begin{equation}
\frac{\rho_\Lambda}{\rho_P} = \frac{7}{2} \times (7\pi)^{-92}
\end{equation}

El exponente 92 viene de:
\begin{itemize}
    \item 46 dimensiones efectivas en teoría de cuerdas/M
    \item Cada una contribuye 2 inversiones
    \item Total: $46 \times 2 = 92$
\end{itemize}


% ============================================================================
\section{Lo que Klein Permite y No Permite}
% ============================================================================

\begin{tcolorbox}[colback=green!10, colframe=green!60!black, title=\textbf{PERMITIDO}]
$\checkmark$ \textbf{Modos con paridad par} bajo TODAS las inversiones de Klein
\begin{itemize}
    \item Aparecen con amplitud completa
    \item Ejemplos: masa del protón, masa del Higgs
\end{itemize}

$\checkmark$ \textbf{Modos mixtos} donde componentes pares dominan
\begin{itemize}
    \item Aparecen con pequeñas correcciones
    \item Ejemplos: velocidad de la luz, constante de estructura fina
\end{itemize}

$\checkmark$ \textbf{Modos impares} en números pequeños de dimensiones
\begin{itemize}
    \item Suprimidos pero detectables
    \item Ejemplos: masas de neutrinos, asimetría materia-antimateria
\end{itemize}
\end{tcolorbox}

\begin{tcolorbox}[colback=yellow!10, colframe=yellow!60!black, title=\textbf{FUERTEMENTE SUPRIMIDO}]
$\bigtriangleup$ \textbf{Modos que violan muchas simetrías} simultáneamente
\begin{itemize}
    \item Suprimidos exponencialmente por $(7\pi)^{-n}$
    \item Ejemplo: constante cosmológica ($n=92$)
\end{itemize}

$\bigtriangleup$ \textbf{Modos que requieren coherencia} en todas las dimensiones
\begin{itemize}
    \item Extremadamente raros
    \item Ejemplo: correlaciones cosmológicas de gran escala
\end{itemize}
\end{tcolorbox}

\begin{tcolorbox}[colback=red!10, colframe=red!60!black, title=\textbf{NO PERMITIDO}]
$\times$ \textbf{Modos puramente impares} en TODAS las dimensiones
\begin{itemize}
    \item Cancelación exacta por interferencia destructiva
    \item No existen en el espectro físico
\end{itemize}

$\times$ \textbf{Estados que violen la condición de consistencia} de Klein
\begin{itemize}
    \item $f(x,-y) \neq \pm f(x,y)$ para ningún signo
    \item Matemáticamente imposibles en esta topología
\end{itemize}

$\times$ \textbf{Monopolos magnéticos aislados}
\begin{itemize}
    \item La topología Klein no permite cargas magnéticas solitarias
    \item Consistente con la no observación experimental
\end{itemize}
\end{tcolorbox}


% ============================================================================
\section{Regla Predictiva: ¿Cómo Determinar $n$?}
% ============================================================================

Para predecir el exponente $n$ de supresión de una cantidad física:

\textbf{1. IDENTIFICAR VIOLACIONES DE SIMETRÍA:}

\begin{center}
\begin{tabular}{lc}
\toprule
\textbf{Simetría violada} & \textbf{Contribuye} \\
\midrule
Paridad espacial (P) & $+1$ \\
Conjugación de carga (C) & $+1$ \\
Reversión temporal (T) & $+1$ \\
Helicidad/Quiralidad & $+1$ \\
Sabor (mezcla de generaciones) & $+1$ \\
Color (QCD) & $+1$ \\
Número bariónico/leptónico & $+1$ \\
\bottomrule
\end{tabular}
\end{center}

\textbf{2. CONTAR DIMENSIONES INVOLUCRADAS:}

Cada dimensión Klein compacta: $+1$ por inversión requerida

\textbf{3. CALCULAR SUPRESIÓN:}

\begin{equation}
\text{Supresión} = (7\pi)^{-n} \quad \text{donde } n = \Sigma(\text{violaciones})
\end{equation}


% ============================================================================
\section{Resumen del Capítulo}
% ============================================================================

\begin{tcolorbox}[colback=kleingray, colframe=kleinblue, title=\textbf{Resumen del Capítulo 10}]

\textbf{IDEAS CENTRALES:}

\begin{enumerate}
    \item La topología Klein impone \textbf{reglas de selección} sobre los modos físicos.

    \item \textbf{Modos pares} (palíndromos topológicos) están permitidos.

    \item \textbf{Modos impares} están suprimidos por factores $(7\pi)^{-n}$.

    \item El exponente $n$ = número de inversiones/violaciones de simetría.

    \item Esto explica por qué algunas cantidades son enormes y otras diminutas.
\end{enumerate}

\vspace{0.3cm}
\textbf{EXPONENTES CLAVE:}

\begin{center}
\begin{tabular}{cl}
$n=2$ & velocidad de la luz, $\alpha$ \\
$n=5$ & masa del neutrino \\
$n=7$ & asimetría bariónica \\
$n=14$ & jerarquía gravitacional \\
$n=24$ & temperatura CMB \\
$n=45$ & edad del universo \\
$n=92$ & constante cosmológica \\
\end{tabular}
\end{center}

\end{tcolorbox}


% ============================================================================
\section*{Ejercicios}
% ============================================================================

\textbf{NIVEL BÁSICO:}
\begin{enumerate}
    \item Explica con tus palabras por qué una cinta de Möbius solo permite
    funciones ``palíndromo''.

    \item Si un modo tiene 3 inversiones, ¿cuál es su factor de supresión?

    \item ¿Por qué la constante cosmológica está tan suprimida?
\end{enumerate}

\textbf{NIVEL INTERMEDIO:}
\begin{enumerate}
\setcounter{enumi}{3}
    \item Calcula el factor de supresión para un proceso que viola C, P y T
    pero no B ni L.

    \item Usando la regla de conteo, predice el orden de magnitud de la
    probabilidad de desintegración del protón.

    \item ¿Por qué hay exactamente 3 generaciones de fermiones según Klein?
\end{enumerate}

\textbf{NIVEL AVANZADO:}
\begin{enumerate}
\setcounter{enumi}{6}
    \item Deriva la condición de consistencia para campos espinoriales
    en la botella de Klein.

    \item Explica por qué SU(5) tiene dimensión 24 y cómo esto se relaciona
    con el exponente en la predicción del CMB.

    \item Propón un experimento para detectar un modo con $n=10$ inversiones.
    ¿Qué precisión necesitarías?
\end{enumerate}
